% ============================================================================
%  Resumen y Abstract
% ============================================================================
\chapter*{Resumen}
\addcontentsline{toc}{chapter}{Resumen}
\markboth{Resumen}{Resumen}

En los ámbitos industriales, el máximo exponente en simplicidad, desempeño y robustez ha sido el
controlador PID, que abarca casi todos los lazos de control. Por eso se busca un nuevo paradigma en
el que se cambie por controladores más avanzados teóricamente, que ofrecen mejores prestaciones, como
los de rechazo activo de perturbaciones (ADRC) o su equivalencia en control
proporcional-integral generalizado (GPI). En ellos cambia el enfoque de sintonización por uno donde
se coloca la dinámica deseada del error de seguimiento y se rechazan las diferentes perturbaciones
identificadas. Al estar basados en el error de seguimiento, se busca que sean lo más simples posible
para que puedan llegar a sustituir al PID casi \emph{plug and play}.

Como caso de estudio se toma el lazo de control de oxígeno disuelto en la etapa de lodos activados
de una planta de tratamiento de aguas residuales, donde se presentan diferentes perturbaciones, como
los caudales de entrada, los mismos actuadores con sus zonas muertas y algunas endógenas, como las
cargas orgánicas, que afectan los niveles de oxígeno en el tanque. A partir de 22 meses de datos de
operación se identificó la planta y se caracterizó la perturbación, y con eso se diseñaron un ADRC
basado en observador de estado extendido (ESO) y un GPI, que resultan ser el mismo lazo escrito de
dos maneras.

Se presentan también los aspectos prácticos de implementación, como el manejo de la saturación, el
acondicionamiento de la referencia, el manejo del ruido y del retardo, junto con una realización en
tiempo discreto para equipos digitales, tanto para el ESO como para el GPI, además de una metodología
práctica para llevar estos controladores a entornos industriales, con bloques de función en texto
estructurado para PLC.

Los controladores se evalúan por medio de simulación con la planta identificada, sin retardo y con un
retardo de 50~s, frente a un PID ISA con la misma sensibilidad máxima y los mismos cambios de
referencia. Se comparan el PID y el ADRC con métricas clásicas de desempeño y de robustez, como el IAE,
el ITAE, el sobrepico y los márgenes de fase y de ganancia, con el fin de validar en qué se destaca
cada arquitectura.

\vspace{0.4cm}
\noindent
\textbf{Palabras clave:} rechazo activo de perturbaciones, observador de estado
extendido, control proporcional-integral generalizado, oxígeno disuelto, lodos
activados, identificación a partir de datos de operación, anti-windup,
PLC, texto estructurado.

\chapter*{Abstract}
\addcontentsline{toc}{chapter}{Abstract}
\markboth{Abstract}{Abstract}

\noindent
\textbf{Active Disturbance Rejection Control for Dissolved Oxygen Regulation in
the Activated Sludge Stage of a Wastewater Treatment Plant}

\vspace{0.3cm}
\noindent
In industry, the PID controller has been the leading choice for simplicity, performance and
robustness, and it covers almost every control loop. This work therefore looks at replacing it with
theoretically more advanced controllers that offer better capabilities, such as active disturbance
rejection control (ADRC) or its equivalent, generalized proportional-integral (GPI) control. These
controllers change the tuning approach: the desired tracking error dynamics are imposed and the
identified disturbances are rejected. Since they work on the tracking error, the aim is to keep them
simple enough to replace a PID almost as plug and play.

The case study is the dissolved oxygen control loop in the activated sludge stage of a wastewater
treatment plant, which is affected by inlet flows, by the actuators themselves and their dead zones,
and by endogenous disturbances such as organic loads. From 22 months of operating data, the plant was
identified and its disturbance characterized, and an ADRC based on an extended state observer (ESO)
and a GPI were designed, which turn out to be the same loop written in two ways.

Practical implementation issues are also addressed: saturation handling, reference conditioning,
noise and dead time, together with a discrete-time realization for digital equipment, for both the
ESO and the GPI, and a practical methodology for taking these controllers to industrial environments,
with function blocks in structured text for PLCs.

The controllers are evaluated in simulation with the identified plant, without dead time and with a
50~s dead time, against an ISA PID with the same maximum sensitivity and the same setpoint changes.
The PID and the ADRC are compared with classical performance and robustness metrics, such as the IAE,
the ITAE, the overshoot and the phase and gain margins, in order to establish where each architecture
stands out.

\vspace{0.4cm}
\noindent
\textbf{Keywords:} active disturbance rejection control, extended state
observer, generalized proportional-integral control, dissolved oxygen,
activated sludge, identification from operating data, anti-windup, PLC,
structured text.
